%=========== ΛΥΣΗ - 71ο ΘΕΜΑ Γ ===========
\begin{thema}{Γ}
\begin{erwthma}
\item Από τη σχέση
\[
f^2(x)-2xf(x)=1
\]
παίρνουμε
\[
f^2(x)-2xf(x)-1=0
\]
Θεωρώντας την ως δευτεροβάθμια εξίσωση ως προς $f(x)$, έχουμε
\[
f(x)=\frac{2x\pm\sqrt{4x^2+4}}{2}=x\pm\sqrt{x^2+1}
\]
Άρα, για κάθε $x\in\mathbb{R}$, ισχύει
\[
f(x)=x+\sqrt{x^2+1}
\quad \text{ή}\quad
f(x)=x-\sqrt{x^2+1}
\]

Οι δύο συναρτήσεις
\[
g_1(x)=x+\sqrt{x^2+1}
\quad \text{και}\quad
g_2(x)=x-\sqrt{x^2+1}
\]
δεν παίρνουν ποτέ την ίδια τιμή, αφού
\[
g_1(x)-g_2(x)=2\sqrt{x^2+1}>0
\]
Επειδή η $f$ είναι συνεχής στο $\mathbb{R}$, δεν μπορεί να αλλάζει κλάδο, αφού τότε, λόγω συνέχειας, θα έπρεπε σε κάποιο σημείο να ισχύει $g_1(x)=g_2(x)$.

Από την αρχική συνθήκη $f(0)=1$ έχουμε
\[
g_1(0)=1\quad \text{και}\quad g_2(0)=-1
\]
Άρα η $f$ ακολουθεί τον πρώτο κλάδο, δηλαδή
\[
f(x)=x+\sqrt{x^2+1}\ ,\quad x\in\mathbb{R}
\]

\item Για κάθε $x\in\mathbb{R}$ έχουμε
\[
f(x)=x+\sqrt{x^2+1}
\]
και
\[
f(-x)=-x+\sqrt{x^2+1}
\]
Άρα
\begin{align*}
f(x)f(-x)
&=\left(x+\sqrt{x^2+1}\right)\left(-x+\sqrt{x^2+1}\right)\\
&=(x^2+1)-x^2\\
&=1
\end{align*}

Επομένως $f(x)\neq 0$ για κάθε $x\in\mathbb{R}$.
Επειδή η $f$ είναι συνεχής στο $\mathbb{R}$ και $f(0)=1>0$, η $f$ δεν μπορεί να αλλάξει πρόσημο.
Άρα
\[
f(x)>0\ ,\quad \text{για κάθε }x\in\mathbb{R}
\]

\item Από τον τύπο
\[
f(x)=x+\sqrt{x^2+1}
\]
η $f$ είναι παραγωγίσιμη στο $\mathbb{R}$ και
\begin{align*}
f'(x)
&=1+\frac{x}{\sqrt{x^2+1}}\\
&=\frac{\sqrt{x^2+1}+x}{\sqrt{x^2+1}}\\
&=\frac{f(x)}{\sqrt{x^2+1}}
\end{align*}
Από το προηγούμενο ερώτημα είναι $f(x)>0$ για κάθε $x\in\mathbb{R}$, ενώ
\[
\sqrt{x^2+1}>0
\]
Άρα
\[
f'(x)>0\ ,\quad \text{για κάθε }x\in\mathbb{R}
\]
Επομένως η $f$ είναι γνησίως αύξουσα στο $\mathbb{R}$.

\item Η $f$ είναι συνεχής στο $[x,x+1]$ και παραγωγίσιμη στο $(x,x+1)$.
Σύμφωνα με το Θεώρημα Μέσης Τιμής, υπάρχει $\xi\in(x,x+1)$ τέτοιο ώστε
\[
f(x+1)-f(x)=f'(\xi)(x+1-x)=f'(\xi)
\]

Υπολογίζουμε τη δεύτερη παράγωγο:
\begin{align*}
f''(x)
&=\left(1+\frac{x}{\sqrt{x^2+1}}\right)'\\
&=\left(\frac{x}{\sqrt{x^2+1}}\right)'\\
&=\frac{\sqrt{x^2+1}-x\cdot\frac{x}{\sqrt{x^2+1}}}{x^2+1}\\
&=\frac{x^2+1-x^2}{(x^2+1)\sqrt{x^2+1}}\\
&=\frac{1}{(x^2+1)^{\frac{3}{2}}}
\end{align*}
Άρα
\[
f''(x)>0\ ,\quad \text{για κάθε }x\in\mathbb{R}
\]
οπότε η $f'$ είναι γνησίως αύξουσα στο $\mathbb{R}$.

Επειδή
\[
\xi<x+1
\]
έχουμε
\[
f'(\xi)<f'(x+1)
\]
Άρα
\[
f(x+1)-f(x)<f'(x+1)
\]
για κάθε $x\in\mathbb{R}$.
\end{erwthma}
\end{thema}
%=========== ΤΕΛΟΣ ΛΥΣΗΣ - 71ο ΘΕΜΑ Γ ===========
